% !TeX root = main.tex
\lecture{4}{03 Feb 2026 11:00}{Mechanics IV: Perp Axes, Radius of Gyration and Angular Momenta}

\section{Perpendicular Axis Theorem}
This only applies in flat 2D objects, or objects which can be approximated as them, for example a pancake or a coin. Suppose the object sits in the XY plane. The perpendicular axis theorem says that the moment of inertia about the z-axis is:
\[
    I_z = I_x + I_z
\]

For rotation in the x-axis:
\[
    I_x = \sum_i m_i y^2_i
\]

And in the y:
\[
    I_y = \sum_i m_i x^2_i
\]

By Pythag:
\[
    r_i^2 = x^2 + y^2_i
\]
\[
    \implies \sum_i m_i x_i^2 + \sum_i y_i^2 = \sum m_i r_i^2
\]
\[
    I_y + I_x = I_z
\]

\subsection{Spinning a Coin}
Consider spinning a coin vertically on a table. Suppose the coin spins about the y-axis (vertical) and the z-axis is the axis of symmetry. 

By symmetry, $I_x = I_y = I_\text{end}$. By the perpendicular axes theorem, $I_x + I_y = 2 I_\text{end} = I_z$.

Hence:
\[
    2 I_\text{end} = I_z = \frac{MR^2}{2} \implies I_\text{end} = \frac{MR^2}{4} 
\]

\section{Radius of Gyration}
The radius of gyration, $K$ is the effective radius at which we could locate all of the mass of an extended body to get the same moment of inertia about an axis passing through the CoM.

It is defined by:
\[
    I_0 = M K^2 \implies K = \sqrt{\frac{I_0}{M}}
\]

We know that $I_z = \frac{1}{2}MR^2$ and by the definition of the radius of gyration:
\[
    I_z = MK^2
\]

Hence:
\[
    \frac{1}{2}MR^2 = MK^2 \implies K^2 = R^2 / 2
\]

\section{Angular Momentum}
We have defined angular momentum as:
\[
    \vec{L} =\vec{r} \wedge \vec{p}
\]

Angular momentum, just like linear momentum, is conserved. ``Conserved'' means that it is not time dependant, i.e.:
\[
    \frac{d \vec{L}}{dt} = 0
\]

\[
    \frac{d \vec{L}}{dt} = \frac{d \vec{r}}{dt} \times \vec{p} + \vec{r} \times \frac{d \vec{p}}{dt}
\]
\[
    = \vec{v} \times \vec{p} + \vec{r} \times \vec{F}
\]
\[
    = \vec{v} \times m\vec{v} + \vec{r} \times \vec{F}
\]
\[
    = 0 + \vec{r} \times \vec{F}
\]

This second term is torque. There is always a centripetal force for circular motion, so $| \vec{F} |  = 0$ does not hold. However, as $\vec{r}$ and $\vec{F}$ are in the same direction, this term does become zero. Accordingly, angular momentum is conserved under radially inwards (central) forces.

And since $\vec{r}$ is defined around a specific origin point, angular momentum is defined about a point in space and varies from point to point.


